\documentclass[11pt]{amsart}

\usepackage[margin=1.15in]{geometry}
\usepackage{amsmath,amssymb,amsthm,mathtools}
\usepackage{enumitem}
\usepackage{microtype}
\usepackage{hyperref}
\usepackage[nameinlink,capitalize]{cleveref}

\hypersetup{
  colorlinks=true,
  linkcolor=blue,
  citecolor=blue,
  urlcolor=blue
}

\theoremstyle{plain}
\newtheorem{theorem}{Theorem}[section]
\newtheorem{proposition}[theorem]{Proposition}
\newtheorem{lemma}[theorem]{Lemma}
\newtheorem{corollary}[theorem]{Corollary}

\theoremstyle{definition}
\newtheorem{definition}[theorem]{Definition}
\newtheorem{hypothesis}[theorem]{Hypothesis}

\theoremstyle{remark}
\newtheorem{remark}[theorem]{Remark}

\numberwithin{equation}{section}

\newcommand{\R}{\mathbb R}
\newcommand{\N}{\mathbb N}
\newcommand{\eps}{\varepsilon}
\newcommand{\loc}{\mathrm{loc}}
\newcommand{\supp}{\operatorname{supp}}
\newcommand{\divergence}{\operatorname{div}}
\newcommand{\calH}{\mathcal H}

\title[Uniformly tight ancient dynamics]
{Uniformly Tight Ancient Navier--Stokes Dynamics and Conditional Equilibrium Rigidity}

\author{Guillem Duran Ballester}
\address{Fragile Tech}
\email{guillem@fragile.tech}
\date{}

\subjclass[2020]{35Q30, 35B44, 35B41, 37A25, 76D05}
\keywords{Navier--Stokes equations, ancient solutions, self-similar variables, invariant measures, Lyapunov functionals, Liouville theorems}

\begin{document}

\begin{abstract}
We study bounded ancient solutions of the three-dimensional Navier--Stokes
equations in backward self-similar variables under uniform \(L^3\)-tightness.
We first establish the analytic compactness framework: bounded mild ancient
solutions are smooth, locally compact under time translation, and the
uniform \(L^\infty_\tau L^3_y\) and tightness bounds are stable under locally
smooth limits.  We then prove a minimal-element reduction for closed normalized
uniformly tight collections: any nonempty such collection contains a nonzero
solution of minimal \(L^\infty(\R^3\times\R)\) size.

For the time-translation hull of a bounded ancient solution we construct compact
invariant hulls and invariant probability measures by Krylov--Bogolyubov
averaging.  These measures satisfy a stationary statistical form of the
self-similar Navier--Stokes equation.  Passing to the barycenter produces an
explicit quadratic covariance term, isolating the obstruction to obtaining a
single stationary self-similar profile.

Finally, we formulate a strict compact-hull Lyapunov hypothesis.  Under this
hypothesis, invariant measures are supported on stationary trajectories.  Hence
a minimal compact uniformly \(L^3\)-bounded hull consists of a single stationary
profile, which is zero by the stationary self-similar theorem of
Ne\v{c}as--R\r{u}\v{z}i\v{c}ka--\v{S}ver\'ak.  As a consequence, the uniformly
\(L^3\)-tight branch of the Type I ancient-profile reduction is excluded
provided the corresponding nonzero compact-recurrence and generated Lyapunov
properties hold.
\end{abstract}

\maketitle

\section{Introduction and main results}\label{sec:introduction}

\subsection{The uniformly tight ancient branch}

We consider the three-dimensional incompressible Navier--Stokes equations in
backward self-similar variables.  The centered velocity \(V\) and pressure \(P\)
solve
\begin{equation}\label{eq:centered}
   \partial_\tau V+(V\cdot\nabla)V+\nabla P
   =
   \Delta V-\frac12(V+y\cdot\nabla V),
   \qquad \divergence V=0
\end{equation}
on \(\R^3\times\R\).  In the companion Type I reduction paper
\cite{DuranPaperI}, admissible Type I singularities are reduced to nonzero
normalized ancient solutions of \eqref{eq:centered}.  The present paper starts
after that reduction and studies the branch in which
\begin{equation}\label{eq:L3-bound}
   \sup_{\tau\in\R}\|V(\tau)\|_{L^3(\R^3)}<\infty
\end{equation}
and
\begin{equation}\label{eq:L3-tightness}
   \lim_{R\to\infty}
   \sup_{\tau\in\R}\int_{|y|>R}|V(y,\tau)|^3\,dy=0 .
\end{equation}
We call this the uniformly \(L^3\)-tight ancient branch.

This profile condition is distinct from the source-level condition
\[
   u\in L^\infty(0,T;L^3(\R^3)).
\]
For finite-energy suitable weak solutions, the source-level assumption
\(u\in L^\infty_tL^3_x\) near a putative singular time is incompatible with
singularity by the theorem of Escauriaza--Seregin--\v{S}ver\'ak
\cite{ESS2003}.  Uniform \(L^3\)-tightness in this paper is a condition on
centered ancient profiles after blow-up.

\subsection{Minimal elements and compact hulls}

The first reduction is variational:
\[
   \hbox{nonempty tight normalized class}
   \quad\Longrightarrow\quad
   \hbox{\(L^\infty\)-minimal ancient element}.
\]
The normalization is a fixed local velocity nonvanishing condition.  It
prevents minimizing sequences from losing all mass by time translation.
Once a minimal element \(V_*\) is selected, its time-translation hull is
\[
   \calH(V_*)
   =
   \overline{\{\Theta_sV_*:\ s\in\R\}}^{C^\infty_{\loc}},
   \qquad
   (\Theta_sV)(y,\tau)=V(y,\tau+s).
\]
The hull is compact because bounded mild ancient solutions are locally smooth
with estimates depending only on their common \(L^\infty\) bound.

\subsection{Invariant measures and covariance}

Every compact invariant hull \(K\) carries invariant probability measures.
Indeed, Krylov--Bogolyubov averages
\[
   \mu_T=\frac1T\int_0^T\delta_{\Theta_sW_0}\,ds
\]
have weak subsequential limits.  Such invariant measures satisfy a stationary
statistical form of \eqref{eq:centered}.  If
\[
   \overline W(y)=\int_K W(y,0)\,d\mu(W),
\]
then the averaged nonlinearity decomposes as
\[
   \int_K W_iW_j\,d\mu
   =
   \overline W_i\,\overline W_j+C_{ij}.
\]
Thus the barycenter solves the stationary self-similar equation only up to the
explicit covariance defect \(C\).  The invariant-measure argument alone does
not prove that a recurrent compact hull contains a single stationary profile.

\subsection{The strict Lyapunov hypothesis}

The present paper therefore isolates the exact conditional input needed to pass
from compact recurrent dynamics to equilibria.  It does not construct a
universal Lyapunov functional for all compact ancient Navier--Stokes hulls.
Instead, it proves that if such a strict Lyapunov functional exists on the
compact hull generated by a minimal tight element, then the tight branch is
excluded.

The logical status of the paper is summarized as follows.

\begin{center}
\small
\begin{tabular}{p{0.58\linewidth}p{0.25\linewidth}}
\hline
Step & Status in this paper\\
\hline
bounded ancient smoothing & proved\\
local smooth compactness & proved\\
closure of uniform \(L^3\)-tightness & proved\\
minimal \(L^\infty\) element & proved\\
compact trajectory hull & proved\\
invariant measures & proved\\
stationary statistical identity & proved\\
covariance obstruction & identified\\
strict Lyapunov functional & assumed\\
nonzero compact recurrence of minimal element & assumed\\
compact tight rigidity under Lyapunov & proved\\
tight-class input for the companion paper & conditional corollary\\
\hline
\end{tabular}
\end{center}

Thus the paper does not prove an unconditional uniformly tight Liouville
theorem.  It proves that the uniformly tight branch is empty once the two
compact-dynamical inputs \(\mathrm{RC}\) and \(\mathrm{LY}\) formulated in
\Cref{def:RC,def:LY} are available.

\subsection{Main results}

The first headline theorem is the analytic compactness and closure statement.

\begin{theorem}[Compactness and closure of uniformly tight ancient solutions]
\label{thm:tight-compactness-intro}
Let \(\{V_n\}\) be a sequence of bounded mild ancient solutions of
\eqref{eq:centered}.  Assume
\[
   \sup_n\|V_n\|_{L^\infty(\R^3\times\R)}<\infty .
\]
Then a subsequence converges locally smoothly to a bounded mild ancient
solution \(V\).

If in addition there are \(L<\infty\) and a modulus
\(\omega:[1,\infty)\to[0,\infty)\), \(\omega(R)\to0\), such that
\[
   \sup_n\sup_{\tau\in\R}\|V_n(\tau)\|_{L^3(\R^3)}\le L
\]
and
\[
   \sup_n\sup_{\tau\in\R}
   \int_{|y|>R}|V_n(y,\tau)|^3\,dy\le \omega(R)
   \qquad (R\ge1),
\]
then
\[
   \sup_{\tau\in\R}\|V(\tau)\|_{L^3(\R^3)}\le L,
   \qquad
   \sup_{\tau\in\R}\int_{|y|>R}|V(y,\tau)|^3\,dy\le\omega(R),
\]
and in particular \(V\) is uniformly \(L^3\)-tight.
\end{theorem}

Compatible pressure representatives also converge locally weakly modulo
time-dependent gauges under the additional \(L^\infty_\tau L^3_y\) assumption;
see \Cref{lem:pressure-reconstruction}.

The second headline theorem is the minimal-element reduction.

\begin{theorem}[Minimal element in a normalized uniformly tight class]
\label{thm:minimal-tight-intro}
Let \(A\) be a nonempty collection of bounded mild ancient solutions satisfying
a common \(L^\infty\) bound, a common \(L^\infty_\tau L^3_y\) bound, a common
uniform \(L^3\)-tightness modulus, and a fixed local nonvanishing
normalization
\[
   \sup_{\tau\in\R}\int_{B_{R_0}}|V(y,\tau)|^3\,dy\ge \eta_0>0.
\]
Assume \(A\) is invariant under time translation and closed under locally
smooth limits which retain the normalization.  Then \(A\) contains an element
\(V_*\) minimizing
\[
   \|V\|_{L^\infty(\R^3\times\R)}
\]
among all \(V\in A\).  In particular \(V_*\not\equiv0\).
\end{theorem}

The third headline theorem is the conditional tight-branch Liouville
criterion.

\begin{theorem}[Conditional uniformly tight Liouville criterion]
\label{thm:conditional-tight-liouville-intro}
Let \(A\) be a normalized uniformly \(L^3\)-tight collection as in
\Cref{thm:minimal-tight-intro}.  If \(A\) satisfies the recurrence property
\(\mathrm{RC}(A)\) and the generated Lyapunov property \(\mathrm{LY}(A)\) of
\Cref{def:RC,def:LY}, then
\[
   A=\varnothing .
\]
\end{theorem}

The proof of the last theorem follows the dependency diagram
\[
\begin{array}{c}
\text{uniformly \(L^3\)-tight normalized class }A\\[2mm]
\Downarrow\\[2mm]
\text{\(L^\infty\)-minimal element }V_*\\[2mm]
\Downarrow\\[2mm]
\text{compact trajectory hull }\calH(V_*)\\[2mm]
\Downarrow\\[2mm]
\text{nonzero compact minimal invariant hull }K\\[2mm]
\Downarrow\quad\text{strict Lyapunov}\\[2mm]
\text{invariant measures supported on stationary trajectories}\\[2mm]
\Downarrow\\[2mm]
K=\{w\},\quad w\in L^3\\[2mm]
\Downarrow\quad\text{Ne\v{c}as--R\r{u}\v{z}i\v{c}ka--\v{S}ver\'ak}\\[2mm]
w=0\\[2mm]
\Downarrow\\[2mm]
\text{contradiction with nonzero }K .
\end{array}
\]

\subsection{Relation to the companion paper}

The companion paper \cite{DuranPaperI} reduces admissible Type I singularities
to normalized Seregin ancient limits and organizes such limits into small,
stationary \(L^3\), uniformly tight, structured/decay, and residual classes.
The present paper supplies the conditional tight-class input: under the
compact-recurrence and strict Lyapunov properties stated below, the uniformly
\(L^3\)-tight class contains no nonzero normalized ancient limit.

\subsection{Organization of the paper}

\Cref{sec:centered-topology} fixes the centered equation, mild formulation,
time translation, and local smooth topology.  \Cref{sec:smoothing-compactness}
proves smoothing, pressure reconstruction, stability, and compactness.
\Cref{sec:tightness-rigidity} proves closure of uniform \(L^3\)-tightness and
recalls stationary \(L^3\) rigidity.  \Cref{sec:minimal-elements} proves the
minimal-element reduction.  \Cref{sec:compact-hulls} constructs compact
trajectory hulls and minimal invariant subsets.  \Cref{sec:measures-covariance}
constructs invariant measures and identifies the covariance obstruction.
\Cref{sec:lyapunov-equilibrium} proves the consequences of the strict Lyapunov
hypothesis.  \Cref{sec:conditional-liouville} assembles the conditional
tight-branch Liouville criterion and records the interface with the companion
residual-criterion paper.

\section{Centered ancient solutions and the local smooth topology}
\label{sec:centered-topology}

\subsection{The centered equation}

Throughout the paper \(V\) denotes a centered ancient velocity and \(P\) its
pressure.  A generic trajectory in a compact hull is denoted by \(W\), and a
stationary profile by \(w(y)\).  Define
\[
   L V=\Delta V-\frac12 y\cdot\nabla V-\frac12 V .
\]
After applying the Helmholtz--Leray projection \(\mathbb P\), the centered
equation becomes
\begin{equation}\label{eq:projected}
   \partial_\tau V
   =
   LV-\mathbb P\divergence(V\otimes V).
\end{equation}

\subsection{The Ornstein--Uhlenbeck--Stokes semigroup}

Let \(S(a)=e^{aL}\), \(a>0\).  It is given by
\begin{equation}\label{eq:OU-semigroup}
   (S(a)f)(y)
   =
   e^{-a/2}
   \int_{\R^3}
   \frac{\exp\!\left(-\frac{|z-e^{-a/2}y|^2}{4(1-e^{-a})}\right)}
        {[4\pi(1-e^{-a})]^{3/2}}
   f(z)\,dz .
\end{equation}
The factor \(e^{-a/2}\) records the damping caused by the term \(-V/2\).

\begin{lemma}[Centered Stokes kernel estimate]\label{lem:centered-stokes-kernel}
For \(a>0\) and \(\phi\in C_c^\infty(\R^3;\R^3)\),
\begin{equation}\label{eq:stokes-adjoint-estimate}
   \|(S(a)\mathbb P\divergence)^*\phi\|_{L^1(\R^3)}
   \le
   C e^{-a/2}(1-e^{-a})^{-1/2}\|\phi\|_{L^1(\R^3)}.
\end{equation}
Consequently the Duhamel term in \eqref{eq:mild} is well-defined by duality
for bounded \(V\).
\end{lemma}

\begin{proof}
Set \(\lambda=e^{-a/2}\) and \(\theta=1-e^{-a}\).  From
\eqref{eq:OU-semigroup},
\[
   S(a)f(y)=e^{-a/2}(e^{\theta\Delta}f)(\lambda y).
\]
Thus
\[
   S(a)^*\phi(z)
   =
   e^{-a/2}\lambda^{-3}
   (e^{\theta\Delta}\psi)(z),
   \qquad
   \psi(x)=\phi(x/\lambda),
\]
and \(\|\psi\|_{L^1}=\lambda^3\|\phi\|_{L^1}\).  The Leray projection is a
homogeneous Fourier multiplier of degree zero, so it commutes with this
dilation.  The classical heat--Stokes kernel estimate
\[
   \|(e^{\theta\Delta}\mathbb P\divergence)^*\phi\|_{L^1}
   \le C\theta^{-1/2}\|\phi\|_{L^1},
   \qquad \theta>0,
\]
follows from the Oseen kernel bound; see
\cite{FabesJonesRiviere1972,GigaMiyakawa1985,Sohr2001}.  Substituting
\(\theta=1-e^{-a}\) and using
\[
   e^{-a/2}\lambda^{-3}\|\psi\|_{L^1}=e^{-a/2}\|\phi\|_{L^1}
\]
gives \eqref{eq:stokes-adjoint-estimate}.  Since
\((1-e^{-a})^{-1/2}\) is integrable at \(a=0\), the duality pairing in
\eqref{eq:mild} is integrable on finite time intervals for
\(V\in L^\infty\).
\end{proof}

\subsection{Riesz transform convention}

We use the convention that \(R_i\) has Fourier symbol \(i\xi_i/|\xi|\).  Thus
\(R_iR_j\) has symbol \(-\xi_i\xi_j/|\xi|^2\), and
\[
   P=R_iR_j(F_{ij})
   \quad\Longleftrightarrow\quad
   -\Delta P=\partial_i\partial_jF_{ij}.
\]
Pressures are determined only modulo functions of time unless a whole-space
Leray gauge is explicitly chosen.

\subsection{Bounded mild ancient solutions}

\begin{definition}[Bounded mild ancient solution]\label{def:mild}
A vector field \(V:\R^3\times\R\to\R^3\) is a bounded mild ancient solution of
\eqref{eq:centered} if
\[
   V\in L^\infty(\R^3\times\R),\qquad \divergence V=0
\]
in distributions, and for every \(\sigma<\tau\),
\begin{equation}\label{eq:mild}
   V(\tau)
   =
   S(\tau-\sigma)V(\sigma)
   -
   \int_\sigma^\tau
   S(\tau-s)\mathbb P\divergence(V\otimes V)(s)\,ds
\end{equation}
in the weak-* topology of \(L^\infty(\R^3)\).
\end{definition}

The Duhamel term in \eqref{eq:mild} is interpreted by duality.  For
\(\phi\in C_c^\infty(\R^3;\R^3)\),
\[
   \left\langle
      S(\tau-s)\mathbb P\divergence(V\otimes V)(s),\phi
   \right\rangle
   =
   \left\langle
      V\otimes V(s),
      (S(\tau-s)\mathbb P\divergence)^*\phi
   \right\rangle ,
\]
and \Cref{lem:centered-stokes-kernel} makes the scalar integrand absolutely
integrable on \(s\in(\sigma,\tau)\).

\subsection{Time translation}

For \(s\in\R\) define
\[
   (\Theta_sV)(y,\tau)=V(y,\tau+s).
\]
The equation, the mild formulation, the \(L^\infty\) norm, and the
\(L^\infty_\tau L^3_y\) and tail bounds used below are invariant under
\(\Theta_s\).  Indeed, replacing \((\sigma,\tau)\) by
\((\sigma+s,\tau+s)\) in \eqref{eq:mild} and changing variables in the time
integral gives the mild formula for \(\Theta_sV\).

\subsection{The locally smooth topology}

\begin{definition}[Local smooth topology]\label{def:local-smooth-topology}
The locally smooth topology on \(C^\infty_{\loc}(\R^3\times\R;\R^3)\) is
generated by the seminorms
\[
   \|V\|_{m,N}
   =
   \max_{|\alpha|+j\le m}
   \sup_{\substack{|y|\le N\\ |\tau|\le N}}
   |\partial_y^\alpha\partial_\tau^jV(y,\tau)|,
   \qquad m,N\in\N .
\]
One compatible metric is
\[
   d(U,V)
   =
   \sum_{m,N=1}^{\infty}
   2^{-(m+N)}
   \frac{\|U-V\|_{m,N}}{1+\|U-V\|_{m,N}}.
\]
\end{definition}

Thus \(V_n\to V\) locally smoothly if and only if the convergence is in
\(C^m(K)\) for every compact cylinder
\(K\Subset\R^3\times\R\) and every \(m\ge0\).

\section{Smoothing, pressure reconstruction, and compactness}
\label{sec:smoothing-compactness}

\subsection{Local smoothing}

\begin{lemma}[Smoothing of bounded mild solutions]\label{lem:mild-smoothing}
Every bounded mild ancient solution is smooth on \(\R^3\times\R\) and satisfies
the projected equation \eqref{eq:projected} classically.  Moreover, for each
compact cylinder \(K\Subset\R^3\times\R\) and each integer \(m\ge0\),
\[
   \|V\|_{C^m(K)}
   \le
   C\bigl(m,K,\|V\|_{L^\infty(\R^3\times\R)}\bigr).
\]
The same estimate holds uniformly for any family with a common
\(L^\infty\) bound.
\end{lemma}

\begin{proof}
Fix a compact time interval \(J\Subset\R\) containing the time projection of
\(K\), and set \(t=-e^{-\tau}\).  On the corresponding compact interval
\(t(J)\Subset(-\infty,0)\), define
\[
   U(x,t)=(-t)^{-1/2}V(x/\sqrt{-t},-\log(-t)).
\]
The map \((x,t)\mapsto (x/\sqrt{-t},-\log(-t))\) is a smooth diffeomorphism
between \(\R^3\times t(J)\) and \(\R^3\times J\), with all derivatives bounded
on compact subcylinders.

The centered mild formula \eqref{eq:mild} transforms, under this parabolic
change of variables, into the usual projected Navier--Stokes mild formula on
every compact subinterval of \(t(J)\):
\[
   U(t)=e^{(t-s)\Delta}U(s)
   -
   \int_s^t e^{(t-r)\Delta}\mathbb P\divergence(U\otimes U)(r)\,dr,
   \qquad s<t.
\]
This transformation uses only the projected mild equation: the Leray projection
is homogeneous of degree zero and commutes with the parabolic dilation, while
the Ornstein--Uhlenbeck kernel in \eqref{eq:OU-semigroup} becomes the heat
kernel with time parameter \(t-s\).  No pressure formulation is used at this
stage.  The \(L^\infty\) norm of \(U\) on \(\R^3\times t(J)\) is bounded by a
constant depending only on \(J\) and \(\|V\|_{L^\infty}\).

The standard local smoothing theorem for bounded projected mild
Navier--Stokes solutions \cite{Kato1984,GigaMiyakawa1985} gives first local
derivative bounds away from the initial time in any mild representation.  In
the present ancient setting one may choose the initial time strictly before the
compact subcylinder under consideration.  Iterating the differentiated
Duhamel formula and then applying the interior parabolic bootstrap for the
Stokes system \cite{Ladyzhenskaya1968,Sohr2001} gives uniform bounds for all
derivatives of \(U\) on compact subcylinders of \(\R^3\times t(J)\).  The mild
formula may then be differentiated in time, so \(U\) solves the projected
Navier--Stokes equation classically.  Returning to \((y,\tau)\) gives the
asserted estimates for \(V\) and the classical projected equation
\eqref{eq:projected}.  The argument is uniform for families with a common
\(L^\infty\) bound.
\end{proof}

\subsection{Pressure reconstruction}

\begin{lemma}[Local pressure reconstruction]\label{lem:local-pressure}
Let \(V\) be a smooth bounded solution of the projected centered equation
\eqref{eq:projected}.  Then for every ball \(B_R\) and compact interval
\(I\Subset\R\) there exists
\[
   P_R\in C^\infty(B_R\times I),
\]
unique up to an additive function of time, such that
\[
   \partial_\tau V+(V\cdot\nabla)V+\nabla P_R
   =
   \Delta V-\frac12(V+y\cdot\nabla V),
   \qquad \divergence V=0,
\]
on \(B_R\times I\).
\end{lemma}

\begin{proof}
Set
\[
   F=
   \Delta V-\frac12(V+y\cdot\nabla V)
   -\partial_\tau V-(V\cdot\nabla)V .
\]
Since \((V\cdot\nabla)V=\divergence(V\otimes V)\) and \(V\) satisfies
\eqref{eq:projected}, self-adjointness of the Leray projection gives
\[
   \int_I\int_{B_R}F\cdot\Phi\,dy\,d\tau=0
\]
for every \(\Phi\in C_c^\infty(B_R\times I;\R^3)\) satisfying
\(\divergence_y\Phi=0\).  The local de Rham theorem, applied in the spatial
variables with \(\tau\) as a parameter, gives a scalar distribution \(P_R\) on
\(B_R\times I\), unique up to a distribution depending only on time, such that
\(\nabla_yP_R=F\).  Since \(F\) is smooth, all spatial derivatives of \(P_R\)
of positive order are smooth.  Fixing, for instance, the spatial mean of
\(P_R\) over \(B_R\) to be zero for each \(\tau\) selects a representative
which is smooth on \(B_R\times I\).  This proves the local pressure
formulation.
\end{proof}

\begin{lemma}[Leray pressure and local gauges]\label{lem:pressure-reconstruction}
Let \(V\) be a bounded mild ancient solution satisfying
\[
   \sup_{\tau\in\R}\|V(\tau)\|_{L^3(\R^3)}<\infty .
\]
Then the whole-space Leray pressure
\[
   P=R_iR_j(V_iV_j)
\]
belongs to \(L^\infty(\R;L^{3/2}(\R^3))\).  Moreover, if \(V_n\to V\) locally
smoothly, the \(V_n\) have a common \(L^\infty_\tau L^3_y\) bound, and
\[
   P_n=R_iR_j(V_{n,i}V_{n,j}),
\]
then on every compact cylinder \(B_R\times[\tau_1,\tau_2]\) there are bounded
measurable gauges \(a_n(\tau)\) and \(a(\tau)\) such that, after passing to a
subsequence,
\[
   P_n-a_n(\tau)
   \rightharpoonup P-a(\tau)
   \quad\hbox{weakly in }L^{3/2}(B_R\times[\tau_1,\tau_2]).
\]
Equivalently, after replacing the limiting pressure by an equivalent
representative modulo a function of time, the weak limit may be denoted by
\(P\).
\end{lemma}

\begin{proof}
The first assertion follows from the Calderon--Zygmund boundedness of Riesz
transforms on \(L^{3/2}(\R^3)\) \cite{Stein1970}:
\[
   \|P(\tau)\|_{L^{3/2}}
   \le C\|V_i(\tau)V_j(\tau)\|_{L^{3/2}}
   \le C\|V(\tau)\|_{L^3}^2 .
\]
For the local stability statement, choose
\(\chi\in C_c^\infty(B_{4R})\) with \(\chi\equiv1\) on \(B_{3R}\).  Decompose
on each time slice
\[
   P_n=R_iR_j(\chi V_{n,i}V_{n,j})
       +R_iR_j((1-\chi)V_{n,i}V_{n,j}) .
\]
The local term converges strongly in \(L^q(B_{2R}\times[\tau_1,\tau_2])\) for
every finite \(q\).  This is because \(V_n\to V\) smoothly on
\(\overline B_{4R}\times[\tau_1,\tau_2]\), so
\(\chi V_{n,i}V_{n,j}\to\chi V_iV_j\) strongly in every finite \(L^q\), and
the Riesz transforms are Calderon--Zygmund operators on \(L^q\), \(1<q<\infty\).

It remains to control the exterior part modulo constants.  Let \(K_{ij}\) be
the convolution kernel of \(R_iR_j\).  For \(x\in B_R\) and
\(|z|>4R\),
\[
   |K_{ij}(x-z)-K_{ij}(-z)|\le C_R |z|^{-4}.
\]
Define the exterior gauges
\[
   b_n(\tau)
   =
   \int_{\R^3} K_{ij}(-z)(1-\chi(z))V_{n,i}(z,\tau)V_{n,j}(z,\tau)\,dz
\]
and \(b(\tau)\) by the same formula with \(V\) in place of \(V_n\).
The cutoff removes the singularity at the origin, and the integral is
absolutely convergent at infinity because, on the dyadic annulus
\(A_m=\{2^mR<|z|<2^{m+1}R\}\),
\[
   \int_{A_m}|z|^{-3}|V_n(z,\tau)|^2\,dz
   \le
   C(2^mR)^{-3}
   \|V_n(\tau)\|_{L^3(A_m)}^2|A_m|^{1/3}
   \le C L^2(2^mR)^{-2}.
\]
The same dyadic estimate shows that \(b_n\) and \(b\) are bounded uniformly on
\([\tau_1,\tau_2]\).  The maps \(\tau\mapsto b_n(\tau)\) and
\(\tau\mapsto b(\tau)\) are measurable, because they are limits of measurable
truncated integrals.
After subtracting \(b_n(\tau)\), the exterior contribution on \(B_R\) is
\[
   E_n(x,\tau)
   =
   \int_{\R^3}
   [K_{ij}(x-z)-K_{ij}(-z)]
   (1-\chi(z))V_{n,i}(z,\tau)V_{n,j}(z,\tau)\,dz .
\]
The same dyadic estimate, now using \(|z|^{-4}\), gives
\[
   \|E_n(\cdot,\tau)\|_{L^\infty(B_R)}
   \le C(R)L^2
\]
uniformly in \(n\) and \(\tau\).  Splitting the exterior integral into
\(\{4R<|z|<R'\}\) and \(\{|z|>R'\}\), the first part converges by local smooth
convergence, while the second part is bounded uniformly by \(C L^2/R'\).  Letting
first \(n\to\infty\) and then \(R'\to\infty\) identifies the local limit of
\(E_n\) with the corresponding exterior contribution for \(V\).  Therefore,
with \(a_n=b_n\) and \(a=b\), the decompositions of \(P_n-a_n\) and \(P-a\)
have local parts converging strongly and exterior parts converging
distributionally with a uniform local \(L^{3/2}\) bound.  Hence
\[
   P_n-a_n(\tau)\rightharpoonup P-a(\tau)
   \quad\text{weakly in }L^{3/2}(B_R\times[\tau_1,\tau_2]).
\]
\end{proof}

\subsection{Stability under locally smooth limits}

\begin{lemma}[Closure of the mild formulation]\label{lem:mild-limit}
Let \(V_n\) be bounded mild ancient solutions satisfying
\[
   \sup_n\|V_n\|_{L^\infty(\R^3\times\R)}<\infty .
\]
If \(V_n\to V\) locally smoothly, then \(V\) is a bounded mild ancient solution
of \eqref{eq:centered}.
\end{lemma}

\begin{proof}
The common \(L^\infty\) bound passes to \(V\) by pointwise convergence on
compact subsets.  Divergence freeness passes distributionally.  Fix
\(\sigma<\tau\) and test \eqref{eq:mild} for \(V_n\) against
\(\phi\in C_c^\infty(\R^3;\R^3)\).  The endpoint terms are
\[
   \langle V_n(\tau),\phi\rangle,
   \qquad
   \langle V_n(\sigma),S(\tau-\sigma)^*\phi\rangle .
\]
The first converges by local smooth convergence.  For the second, note that
\(S(\tau-\sigma)^*\phi\in L^1(\R^3)\).  Approximate this \(L^1\) function by
a compactly supported smooth function and use the common \(L^\infty\) bound to
control the \(L^1\) approximation error uniformly in \(n\).

For the Duhamel term, write
\[
   \left\langle
      S(\tau-s)\mathbb P\divergence(V_n\otimes V_n)(s),\phi
   \right\rangle
   =
   \left\langle
      V_n(s)\otimes V_n(s),
      (S(\tau-s)\mathbb P\divergence)^*\phi
   \right\rangle .
\]
For fixed \(s\), the adjoint test belongs to \(L^1\); pairing convergence
again follows by compact approximation in \(L^1\), local smooth convergence,
and the common \(L^\infty\) bound on \(V_n\otimes V_n\).  The approximation
error is uniform in \(n\) for this fixed \(s\).  Moreover
\Cref{lem:centered-stokes-kernel} gives the integrable domination
\[
   C\sup_n\|V_n\|_{L^\infty}^2
   e^{-(\tau-s)/2}(1-e^{-(\tau-s)})^{-1/2}\|\phi\|_{L^1},
   \qquad s\in(\sigma,\tau).
\]
Dominated convergence in \(s\) passes the Duhamel integral to the limit.
Thus \(V\) satisfies \eqref{eq:mild}.
\end{proof}

\subsection{Compactness of bounded families}

\begin{proposition}[Compactness of bounded ancient families]
\label{prop:bounded-compactness}
Let \(\{V_n\}\) be bounded mild ancient solutions satisfying
\[
   \sup_n\|V_n\|_{L^\infty(\R^3\times\R)}<\infty .
\]
Then a subsequence converges locally smoothly to a bounded mild ancient
solution \(V\).
\end{proposition}

\begin{proof}
By \Cref{lem:mild-smoothing}, the sequence is bounded in every
\(C^m\)-norm on every compact cylinder.  Arzela--Ascoli on a countable
exhaustion \(B_N\times[-N,N]\), followed by a diagonal extraction over
\((m,N)\), gives a subsequence converging locally smoothly to a smooth limit.
\Cref{lem:mild-limit} shows that the limit is again a bounded mild ancient
solution.
\end{proof}

\section{\texorpdfstring{Uniform \(L^3\)-tightness and stationary \(L^3\) rigidity}{Uniform L3-tightness and stationary L3 rigidity}}
\label{sec:tightness-rigidity}

\subsection{\texorpdfstring{Uniform \(L^3\)-tightness}{Uniform L3-tightness}}

\begin{definition}[Uniform \(L^3\)-tightness]\label{def:L3-tight}
A bounded ancient solution \(V\) is uniformly \(L^3\)-tight if
\[
   \sup_{\tau\in\R}\|V(\tau)\|_{L^3(\R^3)}<\infty
\]
and
\[
   \lim_{R\to\infty}
   \sup_{\tau\in\R}
   \int_{|y|>R}|V(y,\tau)|^3\,dy=0 .
\]
\end{definition}

\subsection{Closedness under locally smooth convergence}

\begin{lemma}[Closure of uniform \(L^3\)-tightness]\label{lem:tight-closure}
Let \(V_n\) be bounded mild ancient solutions satisfying
\[
   \sup_n\sup_{\tau\in\R}\|V_n(\tau)\|_{L^3(\R^3)}\le L<\infty
\]
and assume that they have a common tightness modulus: for some
\(\omega:[1,\infty)\to[0,\infty)\) with \(\omega(R)\to0\),
\[
   \sup_n\sup_{\tau\in\R}
   \int_{|y|>R}|V_n(y,\tau)|^3\,dy\le \omega(R),
   \qquad R\ge1.
\]
If \(V_n\to V\) locally smoothly, then
\[
   \sup_{\tau\in\R}\|V(\tau)\|_{L^3(\R^3)}\le L
\]
and
\[
   \sup_{\tau\in\R}
   \int_{|y|>R}|V(y,\tau)|^3\,dy\le \omega(R),
   \qquad R\ge1.
\]
In particular \(V\) is uniformly \(L^3\)-tight.
\end{lemma}

\begin{proof}
Fix \(\tau\).  On \(B_R\), local smooth convergence gives
\[
   \int_{B_R}|V(y,\tau)|^3\,dy
   =
   \lim_{n\to\infty}\int_{B_R}|V_n(y,\tau)|^3\,dy
   \le L^3 .
\]
Letting \(R\to\infty\) and using monotone convergence gives
\(\|V(\tau)\|_{L^3}\le L\).  Since \(\tau\) was arbitrary, the same bound is
uniform in time.

For the tail bound, fix \(1\le R<R'\).  On the bounded annulus
\(\{R<|y|<R'\}\), dominated convergence gives
\[
   \int_{R<|y|<R'}|V(y,\tau)|^3\,dy
   =
   \lim_{n\to\infty}
   \int_{R<|y|<R'}|V_n(y,\tau)|^3\,dy
   \le \omega(R).
\]
Letting \(R'\to\infty\) and using monotone convergence gives
\[
   \int_{|y|>R}|V(y,\tau)|^3\,dy\le \omega(R).
\]
The estimate is independent of \(\tau\), so taking the supremum in time gives
the asserted tail bound.  Since \(\omega(R)\to0\), \(V\) is uniformly
\(L^3\)-tight.
\end{proof}

\begin{proof}[Proof of \Cref{thm:tight-compactness-intro}]
The compactness assertion is \Cref{prop:bounded-compactness}.  The
\(L^\infty_\tau L^3_y\) and tightness closure assertions follow from
\Cref{lem:tight-closure} with the displayed constants \(L\) and \(\omega\).
\end{proof}

\subsection{Stationary time-translate limits}

The stationary self-similar profile equation used below is
\begin{equation}\label{eq:stationary-profile}
   (w\cdot\nabla)w+\nabla p
   =
   \Delta w-\frac12(w+y\cdot\nabla w),
   \qquad \divergence w=0 .
\end{equation}

\begin{theorem}[Stationary limits under uniform \(L^3\)-tightness]
\label{thm:stationary-limits-tight}
Let \(V\) be uniformly \(L^3\)-tight.  Let \(\tau_n\in\R\), and suppose
\[
   \Theta_{\tau_n}V\to W
\]
locally smoothly.  If \(W\) is stationary, then \(W\equiv0\).
\end{theorem}

\begin{proof}
The sequence \(\Theta_{\tau_n}V\) has the same \(L^3\) bound and tightness
modulus as \(V\).  By \Cref{lem:tight-closure}, the limit \(W\) belongs to
\(L^\infty_\tau L^3_y\).  If \(W\) is stationary, write
\(W(y,\tau)=w(y)\).  Then \(w\in L^3(\R^3)\).

By \Cref{lem:mild-limit}, \(W\) is a bounded mild ancient solution.  The
smoothness from \Cref{lem:mild-smoothing} and the local pressure lemma
\Cref{lem:local-pressure} give smooth local pressures \(p_N\) on \(B_N\) at
time \(0\).  Since
\(\partial_\tau W=0\), each \(p_N\) satisfies
\[
   (w\cdot\nabla)w+\nabla p_N
   =
   \Delta w-\frac12(w+y\cdot\nabla w)
   \quad\text{on }B_N .
\]
On overlaps the gradients of the local pressures agree, so the pressures
differ only by constants.  Adjusting those constants and passing to the
compatible union gives a scalar distribution \(p\) on \(\R^3\) for which
\eqref{eq:stationary-profile} holds in distributions.  The stationary
\(L^3\) rigidity theorem \Cref{thm:NRS} applies and gives \(w=0\).
\end{proof}

\subsection{Stationary self-similar rigidity}

\begin{theorem}[Stationary self-similar rigidity]\label{thm:NRS}
Let \(w\in C^\infty(\R^3;\R^3)\cap L^3(\R^3)\) solve
\eqref{eq:stationary-profile} in distributions for some scalar distribution
\(p\), with \(\divergence w=0\).  Then \(w\equiv0\).
\end{theorem}

\begin{proof}
This is the theorem of Ne\v{c}as, R\r{u}\v{z}i\v{c}ka, and
\v{S}ver\'ak \cite{NRS1996}, applied to the stationary Leray self-similar
system.  The smooth distributional formulation stated here is a special case
of their weak \(L^3(\R^3)\) theorem.
\end{proof}

\section{Normalized collections and minimal ancient solutions}
\label{sec:minimal-elements}

\subsection{Closed normalized collections}

\begin{definition}[Closed normalized collection]\label{def:closed-normalized}
A collection \(A\) of bounded mild ancient solutions is called a closed
normalized collection if:
\begin{enumerate}[label=\textup{(\roman*)}]
\item there is \(M<\infty\) such that
\[
   \|V\|_{L^\infty(\R^3\times\R)}\le M
   \qquad\text{for all }V\in A;
\]
\item there are \(R_0>0\) and \(\eta_0>0\) such that
\begin{equation}\label{eq:local-nonzero}
   \sup_{\tau\in\R}
   \int_{B_{R_0}}|V(y,\tau)|^3\,dy
   \ge \eta_0
   \qquad\text{for all }V\in A;
\end{equation}
\item \(A\) is invariant under time translation:
\[
   V\in A,\ s\in\R
   \quad\Longrightarrow\quad
   \Theta_sV\in A;
\]
\item if \(V_n\in A\), \(V_n\to V\) locally smoothly, and \(V\) satisfies the
same \((R_0,\eta_0)\)-nonvanishing condition \eqref{eq:local-nonzero}, then
\(V\in A\).
\end{enumerate}
\end{definition}

\subsection{Time-centering}

\begin{lemma}[Time-centering]\label{lem:time-centering}
Let \(A\) be a closed normalized collection.  If \(V\in A\), then for every
\(\delta>0\) there exists \(s\in\R\) such that
\[
   \int_{B_{R_0}}|\Theta_sV(y,0)|^3\,dy
   \ge \eta_0-\delta.
\]
\end{lemma}

\begin{proof}
By \eqref{eq:local-nonzero}, the supremum in time of the local \(L^3\) mass is
at least \(\eta_0\).  Choose \(s\) within \(\delta\) of this supremum.  Since
\((\Theta_sV)(y,0)=V(y,s)\), the displayed estimate follows.  Time-translation
invariance keeps \(\Theta_sV\) in \(A\).
\end{proof}

\subsection{Stability of local nonvanishing}

\begin{lemma}[Stability of local nonvanishing]\label{lem:nonzero-stability}
If \(V_n\to V\) locally smoothly and
\[
   \int_{B_{R_0}}|V_n(y,0)|^3\,dy\ge \eta_0-o(1),
\]
then
\[
   \int_{B_{R_0}}|V(y,0)|^3\,dy\ge \eta_0 .
\]
\end{lemma}

\begin{proof}
Local smooth convergence implies uniform convergence on
\(\overline B_{R_0}\times\{0\}\), hence convergence in
\(L^3(B_{R_0})\).  The map \(f\mapsto\int_{B_{R_0}}|f|^3\) is continuous on
\(L^3(B_{R_0})\), and the claim follows by passing to the limit.
\end{proof}

\subsection{Lower semicontinuity of size}

\begin{lemma}[Lower semicontinuity of \(L^\infty\) size]
\label{lem:lsc-size}
If \(V_n\to V\) locally smoothly and
\[
   \sup_n\|V_n\|_{L^\infty(\R^3\times\R)}<\infty,
\]
then
\[
   \|V\|_{L^\infty(\R^3\times\R)}
   \le
   \liminf_{n\to\infty}
   \|V_n\|_{L^\infty(\R^3\times\R)}.
\]
\end{lemma}

\begin{proof}
Let \(L_0\) be the liminf on the right.  If
\(|V(y_0,\tau_0)|>L_0+\eps\), choose a subsequence with
\(\|V_n\|_{L^\infty}\to L_0\).  Local uniform convergence at
\((y_0,\tau_0)\) gives
\(|V_n(y_0,\tau_0)|>L_0+\eps/2\) for all large \(n\), contradicting
\(\|V_n\|_{L^\infty}\to L_0\).  Hence \(|V|\le L_0\) pointwise.  Since the
smooth representative has pointwise supremum equal to its essential
\(L^\infty\) norm, the result follows.
\end{proof}

\subsection{Existence of a minimal element}

\begin{theorem}[Existence of a minimal ancient solution]
\label{thm:minimal-existence}
Let \(A\) be nonempty and closed normalized.  Set
\[
   m=
   \inf_{V\in A}
   \|V\|_{L^\infty(\R^3\times\R)}.
\]
Then \(0<m<\infty\), and there exists \(V_*\in A\) such that
\[
   \|V_*\|_{L^\infty(\R^3\times\R)}=m.
\]
In particular \(V_*\not\equiv0\).
\end{theorem}

\begin{proof}
The common \(L^\infty\) bound gives \(m<\infty\).  The local nonvanishing gives
a positive lower bound:
\[
   \eta_0
   \le
   \sup_{\tau\in\R}\int_{B_{R_0}}|V(y,\tau)|^3\,dy
   \le
   |B_{R_0}|\|V\|_{L^\infty(\R^3\times\R)}^3,
\]
so
\[
   m\ge \left(\eta_0/|B_{R_0}|\right)^{1/3}>0.
\]

Choose \(V_n\in A\) with
\[
   m\le \|V_n\|_{L^\infty(\R^3\times\R)}\le m+\frac1n .
\]
By \Cref{lem:time-centering}, after replacing \(V_n\) by a time translate,
which stays in \(A\) and preserves the \(L^\infty\) norm, we may assume
\[
   \int_{B_{R_0}}|V_n(y,0)|^3\,dy\ge \eta_0-\frac1n .
\]
The common \(L^\infty\) bound and \Cref{prop:bounded-compactness} give a
locally smooth subsequential limit \(V_*\), which is a bounded mild ancient
solution.  By \Cref{lem:nonzero-stability},
\[
   \int_{B_{R_0}}|V_*(y,0)|^3\,dy\ge\eta_0,
\]
so \(V_*\) retains the fixed \((R_0,\eta_0)\)-normalization.  Closedness of
\(A\) then gives \(V_*\in A\).  Finally, \Cref{lem:lsc-size} gives
\[
   \|V_*\|_{L^\infty}
   \le
   \liminf_{n\to\infty}\|V_n\|_{L^\infty}
   =m.
\]
Since \(m\) is the infimum over \(A\) and \(V_*\in A\), equality holds.  The
positive lower bound on \(m\) rules out \(V_*\equiv0\).
\end{proof}

\subsection{The uniformly tight normalized class}

\begin{definition}[Uniformly tight normalized collection]
\label{def:tight-normalized-collection}
Fix \(M,L,R_0,\eta_0>0\) and a modulus
\[
   \omega:[1,\infty)\to[0,\infty),
   \qquad \omega(R)\to0.
\]
A collection \(A\) of bounded mild ancient solutions is called a uniformly
tight normalized collection with data \((M,L,\omega,R_0,\eta_0)\) if:
\begin{enumerate}[label=\textup{(\arabic*)}]
\item for every \(V\in A\),
\[
   \|V\|_{L^\infty(\R^3\times\R)}\le M;
\]
\item for every \(V\in A\),
\[
   \sup_{\tau\in\R}\|V(\tau)\|_{L^3(\R^3)}\le L;
\]
\item for every \(V\in A\) and every \(R\ge1\),
\[
   \sup_{\tau\in\R}
   \int_{|y|>R}|V(y,\tau)|^3\,dy\le \omega(R);
\]
\item for every \(V\in A\),
\[
   \sup_{\tau\in\R}
   \int_{B_{R_0}}|V(y,\tau)|^3\,dy\ge\eta_0;
\]
\item \(A\) is invariant under time translation;
\item \(A\) is closed under locally smooth limits retaining item \textup{(4)}.
\end{enumerate}
\end{definition}

\begin{proposition}[Uniformly tight normalized collections are closed normalized]
\label{prop:tight-closed-normalized}
Every uniformly tight normalized collection is a closed normalized collection.
\end{proposition}

\begin{proof}
The \(L^\infty\) bound, local nonvanishing, time-translation invariance, and
closedness clause in \Cref{def:closed-normalized} are part of
\Cref{def:tight-normalized-collection}.  The only non-formal point is that
the stated \(L^3\) bound and fixed tightness modulus survive locally smooth
limits.  This follows from \Cref{lem:tight-closure}.
\end{proof}

\begin{proof}[Proof of \Cref{thm:minimal-tight-intro}]
By \Cref{prop:tight-closed-normalized}, the collection is closed normalized.
\Cref{thm:minimal-existence} gives the asserted nonzero minimizer.
\end{proof}

\section{Compact trajectory hulls}
\label{sec:compact-hulls}

\subsection{Trajectory hulls}

\begin{definition}[Trajectory hull]\label{def:trajectory-hull}
For a bounded mild ancient solution \(V\), its trajectory hull is
\[
   \calH(V)
   =
   \overline{\{\Theta_sV:\ s\in\R\}}^{C^\infty_{\loc}}.
\]
\end{definition}

\subsection{Compactness of bounded hulls}

\begin{theorem}[Compact trajectory hull]\label{thm:compact-hull}
If \(V\) is a bounded mild ancient solution, then \(\calH(V)\) is compact in
the locally smooth topology and invariant under every time translation
\(\Theta_s\).  Every \(W\in\calH(V)\) is a bounded mild ancient solution with
\[
   \|W\|_{L^\infty(\R^3\times\R)}
   \le
   \|V\|_{L^\infty(\R^3\times\R)}.
\]
\end{theorem}

\begin{proof}
Every translate \(\Theta_sV\) is a bounded mild ancient solution with the same
\(L^\infty\) bound as \(V\).  Hence every sequence of translates has a
locally smoothly convergent subsequence by \Cref{prop:bounded-compactness}.
If \((W_n)\subset\calH(V)\), choose translates \(\Theta_{s_n}V\) with
\(d(W_n,\Theta_{s_n}V)<1/n\).  A subsequence of \(\Theta_{s_n}V\) converges
locally smoothly, and the corresponding subsequence of \(W_n\) has the same
limit.  Thus \(\calH(V)\) is sequentially compact.  Since the locally smooth
topology is metrizable, it is compact.

The \(L^\infty\) estimate and the mild equation pass to elements of the hull
by local smooth convergence and \Cref{lem:mild-limit}.  For invariance, note
that \(\Theta_t\) maps the orbit of \(V\) onto itself.  If \(W_n\to W\)
locally smoothly, then \(\Theta_tW_n\to\Theta_tW\) locally smoothly because
time translation only shifts the compact time interval on which the
seminorms are evaluated.  Hence \(\Theta_t\calH(V)=\calH(V)\).

\end{proof}

\begin{corollary}[Uniform tightness on hulls]\label{cor:tightness-on-hulls}
If \(V\) is uniformly \(L^3\)-tight with bound \(L\) and modulus \(\omega\),
then every \(W\in\calH(V)\) satisfies
\[
   \sup_{\tau\in\R}\|W(\tau)\|_{L^3(\R^3)}\le L
\]
and
\[
   \sup_{\tau\in\R}\int_{|y|>R}|W(y,\tau)|^3\,dy\le\omega(R),
   \qquad R\ge1.
\]
The same estimates hold uniformly on every compact invariant subset
\(K\subset\calH(V)\).
\end{corollary}

\begin{proof}
Every translate of \(V\) has the same \(L^3\) bound and tightness modulus.
Applying \Cref{lem:tight-closure} to any sequence of translates converging
locally smoothly to \(W\in\calH(V)\) gives the displayed estimates for \(W\).
Taking the supremum over \(W\in K\subset\calH(V)\) gives the final assertion.
\end{proof}

\subsection{Compact invariant sets}

\begin{definition}[Compact invariant set]\label{def:compact-invariant}
A nonempty compact subset \(K\) of the locally smooth space of bounded mild
ancient solutions is invariant if
\[
   \Theta_tK=K
   \qquad\text{for every }t\in\R .
\]
It is minimal if it contains no proper nonempty compact invariant subset.
\end{definition}

\begin{lemma}[Continuity of the translation action]\label{lem:action-continuity}
Let \(K\) be a compact invariant set.  For each \(t\in\R\),
\(\Theta_t:K\to K\) is a homeomorphism, and the action
\[
   \R\times K\ni(t,W)\mapsto \Theta_tW\in K
\]
is jointly continuous.
\end{lemma}

\begin{proof}
For fixed \(t\), convergence \(W_n\to W\) locally smoothly implies
\[
   \|\Theta_tW_n-\Theta_tW\|_{C^m(B_R\times[-T,T])}
   =
   \|W_n-W\|_{C^m(B_R\times[-T+t,T+t])}\to0.
\]
Thus \(\Theta_t\) is continuous, and its inverse is \(\Theta_{-t}\).  If
\(t_n\to t\) and \(W_n\to W\), then on any fixed compact cylinder the
difference
\[
   \Theta_{t_n}W_n-\Theta_tW
\]
is bounded by a term involving \(W_n-W\) on a slightly larger compact cylinder
and a term involving the uniform continuity of finitely many derivatives of
\(W\) on that larger cylinder.  Both terms tend to zero, proving joint
continuity.
\end{proof}

\subsection{Minimal compact invariant subsets}

\begin{proposition}[Minimal compact invariant subsets]\label{prop:minimal-subset}
Let \(K\) be a nonempty compact invariant set.  Then \(K\) contains a
nonempty compact minimal invariant subset.
\end{proposition}

\begin{proof}
Let \(\mathcal F\) be the family of nonempty compact invariant subsets of
\(K\), ordered by reverse inclusion.  For a chain \(\{K_\alpha\}\), the
intersection \(\bigcap_\alpha K_\alpha\) is nonempty by compactness and the
finite intersection property, since the chain is totally ordered by ordinary
inclusion.  The intersection is compact and invariant: invariance follows
from \(\Theta_tK_\alpha=K_\alpha\) for all \(\alpha\) and the fact that
\(\Theta_t\) is a homeomorphism.  Thus every chain has an upper bound in the
reverse-inclusion order.  Zorn's lemma gives a maximal element for that order,
equivalently a minimal compact invariant subset under ordinary inclusion.
\end{proof}

\subsection{The zero trajectory}

\begin{lemma}[The zero trajectory]\label{lem:zero-trajectory}
The singleton \(\{0\}\) is compact and invariant.  If a minimal compact
invariant set contains the zero trajectory, then it equals \(\{0\}\).
\end{lemma}

\begin{proof}
The zero solution satisfies \(\Theta_t0=0\) for all \(t\).  Therefore
\(\{0\}\) is a nonempty compact invariant subset of any invariant set
containing zero.  Minimality forces equality.
\end{proof}

\section{Invariant measures and the covariance obstruction}
\label{sec:measures-covariance}

\subsection{Krylov--Bogolyubov averages}

The following is the standard Krylov--Bogolyubov averaging argument
\cite{KryloffBogoliouboff1937,Billingsley1999}.

\begin{theorem}[Krylov--Bogolyubov averages]\label{thm:KB}
Let \(K\) be a compact invariant set and \(W_0\in K\).  For \(T>0\), define
\[
   \mu_T
   =
   \frac1T\int_0^T\delta_{\Theta_sW_0}\,ds .
\]
Then every sequence \(T_n\to\infty\) has a subsequence such that
\[
   \mu_{T_n}\rightharpoonup\mu
\]
weakly as Borel probability measures on \(K\), and every such limit \(\mu\)
is invariant:
\[
   (\Theta_t)_\#\mu=\mu
   \qquad\text{for all }t\in\R .
\]
If \(K\) is minimal, then every invariant Borel probability measure on \(K\)
has support equal to \(K\).
\end{theorem}

\begin{proof}
The orbit map \(s\mapsto\Theta_sW_0\) is continuous by
\Cref{lem:action-continuity}; thus \(\mu_T\) is the push-forward of normalized
Lebesgue measure on \([0,T]\).  Since \(K\) is compact and metrizable, the
space of Borel probability measures on \(K\) is weakly sequentially compact
\cite{Billingsley1999}.  Hence \(\mu_{T_n}\) has a weakly convergent
subsequence.

For \(F\in C(K)\) and \(t>0\),
\[
   \int_K F(\Theta_tW)\,d\mu_T(W)-\int_KF(W)\,d\mu_T(W)
   =
   \frac1T\int_T^{T+t}F(\Theta_sW_0)\,ds
   -
   \frac1T\int_0^tF(\Theta_sW_0)\,ds .
\]
The absolute value is at most \(2t\|F\|_{C(K)}/T\).  Letting
\(T=T_{n_k}\to\infty\) gives invariance for \(t>0\); negative times follow by
applying the positive-time result to \(F\circ\Theta_t\).

If \(K\) is minimal and \(\mu\) is invariant, then \(\supp\mu\) is nonempty,
closed, and invariant.  Indeed, if \(U\) is a neighborhood of
\(\Theta_tW\), then \(\Theta_{-t}U\) is a neighborhood of \(W\), and invariance
shows \(\mu(U)=\mu(\Theta_{-t}U)>0\) whenever \(W\in\supp\mu\).  Hence
\(\Theta_t(\supp\mu)\subset\supp\mu\).  Applying the same argument with
\(-t\) gives the reverse inclusion, so \(\supp\mu\) is invariant.  It is
nonempty because \(\mu(K)=1\), and it is compact because it is closed in
compact \(K\).  Minimality forces \(\supp\mu=K\).
\end{proof}

\subsection{Support on minimal hulls}

\begin{corollary}[Nonzero support]\label{cor:nonzero-support}
Let \(K\) be a minimal compact invariant set with \(K\ne\{0\}\).  Then zero is
not in \(K\), and every invariant probability measure on \(K\) has support
equal to \(K\).
\end{corollary}

\begin{proof}
If zero were in \(K\), \Cref{lem:zero-trajectory} would give \(K=\{0\}\),
contrary to the assumption.  The support conclusion is the minimal-support part
of \Cref{thm:KB}.
\end{proof}

\subsection{Stationary statistical identity}

\begin{theorem}[Stationary statistical form]\label{thm:statistical-identity}
Let \(K\) be a compact invariant set and \(\mu\) an invariant Borel
probability measure on \(K\).  Then for every solenoidal
\(\phi\in C_c^\infty(\R^3;\R^3)\),
\begin{equation}\label{eq:statistical-identity}
\begin{aligned}
0
=
\int_K
\bigg[
   &\int_{\R^3}
   W(y,0)\cdot
   \left(\Delta\phi+\phi+\frac12 y\cdot\nabla\phi\right)\,dy  \\
   &+
   \int_{\R^3}
   W_i(y,0)W_j(y,0)\partial_j\phi_i(y)\,dy
\bigg]\,d\mu(W).
\end{aligned}
\end{equation}
\end{theorem}

\begin{proof}
Define
\[
   F(W)=\int_{\R^3}W(y,0)\cdot\phi(y)\,dy .
\]
Local smooth convergence makes \(F\) continuous on \(K\): if \(W_n\to W\),
then \(W_n(\cdot,0)\to W(\cdot,0)\) uniformly on the compact support of
\(\phi\).  Invariance gives
\[
   \int_K\frac{F(\Theta_hW)-F(W)}{h}\,d\mu(W)=0
   \qquad (0<|h|\le1).
\]
Choose \(R\) with \(\supp\phi\subset B_R\).  On
\(\overline B_R\times[-1,1]\), compactness of \(K\) in the locally smooth
topology gives
\[
   M_\phi
   =
   \sup_{W\in K}
   \|\partial_\tau W\|_{L^\infty(\overline B_R\times[-1,1])}
   <\infty .
\]
For \(0<|h|\le1\),
\[
   \frac{F(\Theta_hW)-F(W)}{h}
   =
   \int_{\R^3}
   \left(\frac1h\int_0^h\partial_\tau W(y,s)\,ds\right)\cdot\phi(y)\,dy,
\]
and the absolute value is bounded by
\(|B_R|\|\phi\|_{L^\infty}M_\phi\), uniformly in \(W\) and \(h\).  For each
fixed \(W\), smoothness gives uniform convergence of the inner average to
\(\partial_\tau W(y,0)\) on \(\supp\phi\).  Dominated convergence on \(K\)
yields
\[
   \int_K\int_{\R^3}\partial_\tau W(y,0)\cdot\phi(y)\,dy\,d\mu(W)=0 .
\]
For each \(W\), the smooth local equation \eqref{eq:centered} at time zero,
tested against solenoidal \(\phi\), gives
\[
\begin{aligned}
\int_{\R^3}\partial_\tau W\cdot\phi\,dy
=
&\int_{\R^3} W\cdot\Delta\phi\,dy
 +\int_{\R^3} W\cdot\phi\,dy
 +\frac12\int_{\R^3}W\cdot(y\cdot\nabla\phi)\,dy\\
&+\int_{\R^3}W_iW_j\partial_j\phi_i\,dy .
\end{aligned}
\]
The pressure term vanishes because \(\divergence\phi=0\); the drift identity
uses
\[
   \int y_j\partial_jW_i\,\phi_i\,dy
   =
   -3\int W\cdot\phi\,dy-\int W\cdot(y\cdot\nabla\phi)\,dy .
\]
The right-hand side is a continuous bounded function of \(W\in K\).  Averaging
over \(K\) gives \eqref{eq:statistical-identity}.
\end{proof}

\subsection{Barycenter and covariance}

\begin{theorem}[Barycenter and covariance]\label{thm:barycenter-covariance}
Let \(K\) and \(\mu\) be as in \Cref{thm:statistical-identity}.  Define
\[
   \overline W(y)=\int_K W(y,0)\,d\mu(W)
\]
and
\[
   C_{ij}(y)
   =
   \int_K W_i(y,0)W_j(y,0)\,d\mu(W)
   -
   \overline W_i(y)\overline W_j(y).
\]
Then \(\overline W\) and \(C\) are \(C^\infty_{\loc}\), \(\overline W\) is
solenoidal, and for every solenoidal
\(\phi\in C_c^\infty(\R^3;\R^3)\),
\begin{equation}\label{eq:barycenter-covariance}
   \int_{\R^3}
   \overline W\cdot
   \left(\Delta\phi+\phi+\frac12 y\cdot\nabla\phi\right)\,dy
   +
   \int_{\R^3}
   \overline W_i\overline W_j\partial_j\phi_i\,dy
   +
   \int_{\R^3}C_{ij}\partial_j\phi_i\,dy
   =0 .
\end{equation}
In particular, \(\overline W\) is a weak stationary self-similar solution
against solenoidal test fields if and only if
\[
   \int_{\R^3}C_{ij}\partial_j\phi_i\,dy=0
\]
for every solenoidal test field \(\phi\).
\end{theorem}

\begin{proof}
Compactness of \(K\) in the locally smooth topology gives, for every compact
ball and every derivative order, a uniform bound on
\(\partial_y^\alpha W(y,0)\) for \(W\in K\).  Hence one may differentiate
\(\overline W\) under the integral sign.  More explicitly, for every
multiindex \(\alpha\) and every \(y\) in a fixed compact ball,
\[
   \partial_y^\alpha\overline W(y)
   =
   \int_K\partial_y^\alpha W(y,0)\,d\mu(W),
\]
because the difference quotients are dominated by the next derivative
seminorm, uniformly on \(K\).  The same argument applies to the second moment
\[
   M_{ij}(y)=\int_KW_i(y,0)W_j(y,0)\,d\mu(W),
\]
using Leibniz' rule and the uniform bounds on derivatives of \(W\).  This
gives \(C^\infty_{\loc}\) regularity of \(\overline W\), \(M\), and
\(C=M-\overline W\otimes\overline W\).  Since every \(W\) is divergence free,
Fubini against scalar test functions gives \(\divergence\overline W=0\):
for \(\eta\in C_c^\infty(\R^3)\),
\[
   \int_{\R^3}\overline W\cdot\nabla\eta\,dy
   =
   \int_K\int_{\R^3}W(y,0)\cdot\nabla\eta(y)\,dy\,d\mu(W)=0.
\]

For each compactly supported test field \(\phi\), the maps
\[
   W\mapsto\int_{\R^3}W(y,0)\cdot\phi(y)\,dy,
   \qquad
   W\mapsto\int_{\R^3}W_i(y,0)W_j(y,0)\partial_j\phi_i(y)\,dy
\]
are continuous and bounded on \(K\), again by compactness in
\(C^\infty_{\loc}\).  Thus all integrals below are finite, and Fubini's theorem
applies.

Applying \Cref{thm:statistical-identity} and passing the linear term through
the \(K\)-integral gives the first integral in
\eqref{eq:barycenter-covariance}.  Passing the nonlinear term through the
integral gives
\[
   \int_K W_iW_j\,d\mu
   =
   \overline W_i\overline W_j+C_{ij}.
\]
Substitution yields \eqref{eq:barycenter-covariance}.  Comparing this identity
with the pressure-free stationary weak formulation for \(\overline W\) gives
the stated equivalence with vanishing covariance contribution.
\end{proof}

\section{Strict Lyapunov functionals and equilibrium support}
\label{sec:lyapunov-equilibrium}

\subsection{Stationary trajectories}

\begin{definition}[Stationary trajectory]\label{def:stationary-trajectory}
An element \(W\) of a compact invariant hull \(K\) is stationary if
\[
   \Theta_tW=W
   \qquad\text{for every }t\in\R.
\]
Equivalently, \(W(y,\tau)=w(y)\) is independent of \(\tau\).
\end{definition}

\begin{lemma}[Stationary trajectories and stationary profiles]
\label{lem:stationary-profile}
If \(W\) is stationary, then \(W(y,\tau)=w(y)\), and \(w\) solves
\eqref{eq:stationary-profile} in distributions for some scalar distribution
\(p\).
\end{lemma}

\begin{proof}
Stationarity means \(W(y,\tau+t)=W(y,\tau)\) for all \(t\), hence
\(W(y,\tau)=W(y,0)=w(y)\).  Conversely, a time-independent trajectory is fixed
by every \(\Theta_t\).  By \Cref{lem:local-pressure}, for each \(N\) there is
a smooth pressure \(p_N\) on \(B_N\) such that
\[
   (w\cdot\nabla)w+\nabla p_N
   =
   \Delta w-\frac12(w+y\cdot\nabla w)
   \quad\text{on }B_N.
\]
On overlaps, the gradients of two local pressures agree, so their difference
is constant on each connected overlap.  More explicitly, if \(M>N\), then on
\(B_N\) both gradients equal the same smooth vector field
\[
   \Delta w-\frac12(w+y\cdot\nabla w)-(w\cdot\nabla)w .
\]
Thus \(\nabla(p_M-p_N)=0\) on \(B_N\), and \(p_M-p_N\) is constant there.
Choose the additive constants in the pressures inductively so that
\(p_{N+1}=p_N\) on \(B_N\).  If
\(\varphi\in C_c^\infty(\R^3)\), choose \(N\) with
\(\supp\varphi\subset B_N\) and define
\[
   \langle p,\varphi\rangle=\langle p_N,\varphi\rangle .
\]
The compatibility on overlaps makes this independent of \(N\).  Testing the
local identity on a ball containing the support of a vector test field then
shows that \eqref{eq:stationary-profile} holds in \(\mathcal D'(\R^3)\), with
\(\divergence w=0\).
\end{proof}

\subsection{The strict Lyapunov hypothesis}

\begin{hypothesis}[Strict compact-hull Lyapunov functional]
\label{hyp:strict-lyapunov}
Let \(K\) be a compact invariant hull.  There exist continuous functions
\[
   E:K\to\R,
   \qquad
   D:K\to[0,\infty),
\]
such that for every \(W\in K\) and every \(t\ge0\),
\begin{equation}\label{eq:strict-lyapunov}
   E(\Theta_tW)-E(W)
   =
   -\int_0^tD(\Theta_sW)\,ds,
\end{equation}
and
\begin{equation}\label{eq:zero-dissipation}
   D(W)=0
   \quad\Longleftrightarrow\quad
   W \text{ is stationary}.
\end{equation}
\end{hypothesis}

\begin{remark}
The construction of such a Lyapunov pair \((E,D)\) is not asserted in this
paper.  The role of the paper is to prove the consequences of this
hypothesis for uniformly \(L^3\)-tight compact ancient dynamics.
\end{remark}

\subsection{Invariant measures annihilate dissipation}

\begin{lemma}[Invariant measures annihilate Lyapunov dissipation]
\label{lem:zero-dissipation}
Let \(K\) satisfy \Cref{hyp:strict-lyapunov}, and let \(\mu\) be an invariant
Borel probability measure on \(K\).  Then
\[
   \int_KD(W)\,d\mu(W)=0.
\]
\end{lemma}

\begin{proof}
Fix \(t>0\).  Integrating \eqref{eq:strict-lyapunov} with respect to the
invariant measure \(\mu\) gives
\[
   0
   =
   -\int_K\int_0^tD(\Theta_sW)\,ds\,d\mu(W).
\]
The integrand is continuous on the compact product \([0,t]\times K\), so
Fubini applies.  Invariance of \(\mu\) with the continuous test function \(D\)
gives
\[
   \int_K\int_0^tD(\Theta_sW)\,ds\,d\mu(W)
   =
   \int_0^t\int_KD(W)\,d\mu(W)\,ds
   =
   t\int_KD(W)\,d\mu(W).
\]
Since \(t>0\), the claim follows.
\end{proof}

\subsection{Equilibrium support}

\begin{theorem}[Equilibrium support]\label{thm:equilibrium-support}
Let \(K\) be a compact invariant hull satisfying
\Cref{hyp:strict-lyapunov}.  Then every invariant Borel probability measure on
\(K\) is supported on stationary trajectories.
\end{theorem}

\begin{proof}
By \Cref{lem:zero-dissipation},
\[
   \int_KD\,d\mu=0.
\]
Since \(D\ge0\) is continuous, the sets
\[
   K_\eps=\{W\in K:\ D(W)\ge\eps\}
\]
have zero \(\mu\)-measure for every \(\eps>0\).  Therefore
\(\mu(K\setminus\{D=0\})=0\).  By \eqref{eq:zero-dissipation}, \(\{D=0\}\) is
exactly the set of stationary trajectories.  Since \(D\) is continuous, this
stationary set is closed.
\end{proof}

\subsection{Minimal compact hulls are stationary}

\begin{corollary}[Minimal compact hulls are stationary]
\label{cor:minimal-stationary}
Let \(K\) be a minimal compact invariant hull satisfying
\Cref{hyp:strict-lyapunov}.  Then every element of \(K\) is stationary.
Consequently \(K\) consists of a single stationary trajectory.
\end{corollary}

\begin{proof}
By \Cref{thm:KB}, \(K\) carries an invariant Borel probability measure
\(\mu\), and minimality gives \(\supp\mu=K\).  By
\Cref{thm:equilibrium-support}, \(\mu\) is supported on the closed set of
stationary trajectories.  Hence \(K=\supp\mu\) is contained in that set.
If \(W\in K\) is stationary, the singleton \(\{W\}\) is compact and invariant.
Minimality of \(K\) forces \(K=\{W\}\).
\end{proof}

\subsection{Compact tight rigidity}

\begin{theorem}[Rigidity of compact tight minimal dynamics]
\label{thm:compact-tight-rigidity}
Let \(K\) be a minimal compact invariant hull satisfying
\Cref{hyp:strict-lyapunov}.  Assume
\[
   \sup_{W\in K}\sup_{\tau\in\R}
   \|W(\tau)\|_{L^3(\R^3)}<\infty.
\]
Then every element of \(K\) is identically zero.
\end{theorem}

\begin{proof}
By \Cref{cor:minimal-stationary}, \(K=\{W\}\) with
\(W(y,\tau)=w(y)\).  By \Cref{lem:stationary-profile}, \(w\) solves
\eqref{eq:stationary-profile}.  The assumed uniform \(L^3\) bound gives
\(w\in L^3(\R^3)\).  Therefore \Cref{thm:NRS} applies and gives \(w=0\).
\end{proof}

\section{Conditional tight-branch Liouville criterion}
\label{sec:conditional-liouville}

\subsection{Recurrence and generated Lyapunov properties}

From a minimal element \(V_*\), the full trajectory hull \(\calH(V_*)\) is
compact and invariant.  It is not a consequence of compactness alone that this
hull contains a nonzero minimal compact invariant subset; the local limit
dynamics may admit the zero singleton as a minimal invariant subset.  We
therefore isolate the required recurrence input.

\begin{definition}[The recurrence property \(\mathrm{RC}(A)\)]
\label{def:RC}
Let \(A\) be a uniformly tight normalized collection.  We say that \(A\)
satisfies \(\mathrm{RC}(A)\) if every \(L^\infty\)-minimal element
\(V_*\in A\) has a trajectory hull
\[
   \calH(V_*)
\]
containing a compact minimal invariant subset \(K\) with \(K\ne\{0\}\).
\end{definition}

\begin{definition}[The generated Lyapunov property \(\mathrm{LY}(A)\)]
\label{def:LY}
Let \(A\) be a uniformly tight normalized collection satisfying
\(\mathrm{RC}(A)\).  We say that \(A\) satisfies \(\mathrm{LY}(A)\) if every
compact minimal invariant subset \(K\) supplied by \(\mathrm{RC}(A)\) satisfies
the strict compact-hull Lyapunov hypothesis \Cref{hyp:strict-lyapunov}.
\end{definition}

\subsection{The conditional tight Liouville theorem}

\begin{theorem}[No nonempty tight class under compact recurrence and Lyapunov]
\label{thm:conditional-tight-liouville}
No nonempty uniformly tight normalized collection \(A\) can satisfy both
\(\mathrm{RC}(A)\) and \(\mathrm{LY}(A)\).  Equivalently, if \(A\) is a
uniformly tight normalized collection satisfying \(\mathrm{RC}(A)\) and
\(\mathrm{LY}(A)\), then \(A=\varnothing\).
\end{theorem}

\begin{proof}
Suppose \(A\) is nonempty.  By
\Cref{prop:tight-closed-normalized,thm:minimal-existence}, choose an
\(L^\infty\)-minimal element \(V_*\in A\).

By \(\mathrm{RC}(A)\), the trajectory hull \(\calH(V_*)\) contains a compact
minimal invariant subset \(K\) with \(K\ne\{0\}\).  Since \(V_*\) is uniformly
\(L^3\)-tight, \Cref{cor:tightness-on-hulls} gives
\[
   \sup_{W\in K}\sup_{\tau\in\R}\|W(\tau)\|_{L^3(\R^3)}<\infty .
\]
By \(\mathrm{LY}(A)\), \(K\) satisfies the strict Lyapunov hypothesis.  Hence
\Cref{thm:compact-tight-rigidity} implies that every element of \(K\) is
identically zero, so \(K=\{0\}\).  This contradicts \(K\ne\{0\}\).  Therefore
no such nonempty \(A\) exists.
\end{proof}

\begin{proof}[Proof of \Cref{thm:conditional-tight-liouville-intro}]
This is \Cref{thm:conditional-tight-liouville}.
\end{proof}

\subsection{Interface with the companion residual criterion}

\begin{definition}[Tight-class Liouville input]
\label{def:tight-class-input}
The tight-class Liouville input is the assertion that no nonzero normalized
Seregin ancient limit generated by an admissible Type I source is uniformly
\(L^3\)-tight.
\end{definition}

\begin{lemma}[Vanishing velocity removes pressure normalization]
\label{lem:pressure-nonzero-implies-velocity-nonzero}
Let \((V,\Pi)\) be a smooth centered ancient solution of \eqref{eq:centered}.
If \(V\equiv0\), then \(\Pi(y,\tau)=b(\tau)\) for some function of time.
Consequently any pressure-normalized local nonvanishing condition of the form
\[
   \iint_Q\left(|V|^3+|\Pi-b(\tau)|^{3/2}\right)\,dy\,d\tau>0
\]
for every admissible pressure gauge \(b(\tau)\) forces \(V\not\equiv0\).
\end{lemma}

\begin{proof}
If \(V\equiv0\), the centered equation gives \(\nabla\Pi=0\) in distributions.
Thus \(\Pi\) is independent of \(y\) on each time slice, so
\(\Pi(y,\tau)=b(\tau)\).  Choosing this pressure gauge makes both terms in the
displayed pressure-normalized integral vanish on every compact cylinder \(Q\).
\end{proof}

\begin{lemma}[From normalized Seregin nonvanishing to local velocity normalization]
\label{lem:seregin-to-velocity-nonzero}
Let \(V\) be a nonzero smooth bounded normalized Seregin ancient limit.  Then
there exist \(R_0>0\) and \(\eta_0>0\) such that
\[
   \sup_{\tau\in\R}\int_{B_{R_0}}|V(y,\tau)|^3\,dy\ge\eta_0 .
\]
\end{lemma}

\begin{proof}
Since \(V\not\equiv0\), there is \((y_0,\tau_0)\) with
\(V(y_0,\tau_0)\ne0\).  By continuity, \(|V|\ge c>0\) on a small spatial ball
centered at \(y_0\) at time \(\tau_0\).  Choose \(R_0\) large enough to
contain that ball.  Then
\[
   \int_{B_{R_0}}|V(y,\tau_0)|^3\,dy>0.
\]
Taking \(\eta_0\) to be half of this positive value gives the claim.
\end{proof}

\begin{corollary}[Tight-class input for the residual criterion]
\label{cor:tight-input-for-paper-I}
Assume that every uniformly \(L^3\)-tight normalized Seregin collection
generated by admissible Type I sources satisfies \(\mathrm{RC}(A)\) and
\(\mathrm{LY}(A)\).  Then \Cref{def:tight-class-input} holds.
\end{corollary}

\begin{proof}
Suppose, contrary to \Cref{def:tight-class-input}, that the companion
residual-criterion paper produces a nonzero normalized Seregin ancient limit
\(V\) in the uniformly \(L^3\)-tight class.  If one starts instead from a
pressure-normalized nonvanishing condition in that construction, then
\Cref{lem:pressure-nonzero-implies-velocity-nonzero} justifies the conclusion
that \(V\not\equiv0\).  Choose \(M\), \(L\), and a tightness modulus from \(V\),
for instance
\[
   M=\|V\|_{L^\infty(\R^3\times\R)},\qquad
   L=\sup_{\tau\in\R}\|V(\tau)\|_{L^3(\R^3)},
\]
and
\[
   \omega(R)=
   \sup_{\tau\in\R}\int_{|y|>R}|V(y,\tau)|^3\,dy .
\]
By \Cref{lem:seregin-to-velocity-nonzero}, choose \(R_0,\eta_0\) giving local
velocity nonvanishing.  Let
\[
   A_{V,R_0,\eta_0}
   =
   \left\{
      W\in\calH(V):
      \sup_{\tau\in\R}\int_{B_{R_0}}|W(y,\tau)|^3\,dy\ge\eta_0
   \right\}.
\]
This is the normalized compact-hull collection generated by \(V\).  It is
nonempty because \(V\in A_{V,R_0,\eta_0}\).  It is invariant because
\(\calH(V)\) is invariant and the normalization uses a supremum over all
\(\tau\).  If \(W_n\in A_{V,R_0,\eta_0}\), \(W_n\to W\) locally smoothly, and
\(W\) retains the same \((R_0,\eta_0)\)-nonvanishing condition, then
\(W\in\calH(V)\) because \(\calH(V)\) is closed, and hence
\(W\in A_{V,R_0,\eta_0}\).  Thus \(A_{V,R_0,\eta_0}\) has exactly the
closedness property required in \Cref{def:tight-normalized-collection}.

By \Cref{thm:compact-hull,cor:tightness-on-hulls}, every element of
\(\calH(V)\), hence every element of \(A_{V,R_0,\eta_0}\), satisfies the same
\(L^\infty\) bound, \(L^3\) bound, and tightness modulus.  Thus
\(A_{V,R_0,\eta_0}\) is a uniformly tight normalized Seregin collection with
the displayed data.

The assumed recurrence and generated Lyapunov properties apply to this
collection, so \Cref{thm:conditional-tight-liouville} gives
\(A_{V,R_0,\eta_0}=\varnothing\).  This contradicts
\(V\in A_{V,R_0,\eta_0}\).  Hence no such nonzero uniformly tight Seregin
limit exists under the stated assumptions, which is precisely
\Cref{def:tight-class-input}.
\end{proof}

The results above reduce the uniformly \(L^3\)-tight ancient branch to two
compact-dynamical inputs: nonzero compact recurrence of \(L^\infty\)-minimal
elements and the strict Lyapunov equilibrium property on the resulting
minimal compact hulls.  Under these inputs the tight branch is empty, because
the Lyapunov principle forces minimal compact dynamics to be stationary and
the stationary \(L^3\) self-similar theorem forces the stationary profile to
vanish.  Thus, within the present framework, the tight-class Liouville input
used in the companion residual criterion is reduced to the compact recurrence
and strict Lyapunov assertions isolated here.

\begin{thebibliography}{99}

\bibitem{Billingsley1999}
P. Billingsley,
\emph{Convergence of Probability Measures}, second ed.,
Wiley, New York, 1999.

\bibitem{DuranPaperI}
G. Duran Ballester,
\emph{Type I blow-up limits and a residual-class criterion for the
three-dimensional Navier--Stokes equations}, preprint, 2026.

\bibitem{ESS2003}
L. Escauriaza, G. Seregin, and V. \v{S}ver\'ak,
\emph{\(L_{3,\infty}\)-solutions of Navier--Stokes equations and backward
uniqueness},
Russian Math. Surveys 58 (2003), 211--250.

\bibitem{FabesJonesRiviere1972}
E. B. Fabes, B. F. Jones, and N. M. Riviere,
\emph{The initial value problem for the Navier--Stokes equations with data in
\(L^p\)},
Arch. Rational Mech. Anal. 45 (1972), 222--240.

\bibitem{GigaMiyakawa1985}
Y. Giga and T. Miyakawa,
\emph{Solutions in \(L_r\) of the Navier--Stokes initial value problem},
Arch. Rational Mech. Anal. 89 (1985), 267--281.

\bibitem{Kato1984}
T. Kato,
\emph{Strong \(L^p\)-solutions of the Navier--Stokes equation in
\(\R^m\), with applications to weak solutions},
Math. Z. 187 (1984), 471--480.

\bibitem{KryloffBogoliouboff1937}
N. Kryloff and N. Bogoliouboff,
\emph{La theorie generale de la mesure dans son application a l'etude des
systemes dynamiques de la mecanique non lineaire},
Ann. of Math. (2) 38 (1937), 65--113.

\bibitem{Ladyzhenskaya1968}
O. A. Ladyzhenskaya, V. A. Solonnikov, and N. N. Ural'tseva,
\emph{Linear and Quasilinear Equations of Parabolic Type},
American Mathematical Society, Providence, 1968.

\bibitem{NRS1996}
J. Ne\v{c}as, M. R\r{u}\v{z}i\v{c}ka, and V. \v{S}ver\'ak,
\emph{On Leray's self-similar solutions of the Navier--Stokes equations},
Acta Math. 176 (1996), 283--294.

\bibitem{Sohr2001}
H. Sohr,
\emph{The Navier--Stokes Equations: An Elementary Functional Analytic
Approach},
Birkhauser, Basel, 2001.

\bibitem{Stein1970}
E. M. Stein,
\emph{Singular Integrals and Differentiability Properties of Functions},
Princeton University Press, Princeton, 1970.

\end{thebibliography}

\end{document}
